%% it/sections/01_storica.tex
%% ============================================================
%% 01_introduzione_storica.tex
%% Contestualizzazione storica bilingue — circa 10 pagine
%% ============================================================

%% ============================================================
%% VERSIONE ITALIANA
%% ============================================================

\chapter{Il Contesto Storico}
\markboth{Il Contesto Storico}{Il Contesto Storico}

\lettrine[lines=4]{P}{er} comprendere i Vespri Siciliani è necessario
abbandonare la lettura nazionalistica che ne ha dominato la memoria
collettiva — quella del «popolo siciliano che si ribella allo straniero» —
e guardare la Sicilia del 1282 come quello che era: un nodo cruciale in
una rete di conflitti geopolitici che abbracciava l'intero bacino
mediterraneo.

\nota{Nota sul dibattito storiografico}{
  La «cospirazione aragonese» orchestrata da Giovanni da Procida è una
  narrazione consolidata ma storicamente controversa. Storici come
  Steven Runciman (\textit{The Sicilian Vespers}, 1958) e David Abulafia
  la accettano nelle linee generali, mentre altri, tra cui Jean Dunbabin,
  sottolineano quanto la rivolta del 30 marzo fosse in larga misura
  spontanea, alimentata da risentimenti profondi e da un accumulo di
  soprusi quotidiani. \textit{Vespri 1282} sceglie deliberatamente di
  presentare entrambe le letture come livelli sovrapposti della realtà:
  la cospirazione esiste, ma non controlla tutto. È questa tensione che
  rende la storia giocabile.
}

\section{Il Mediterraneo nel XIII secolo: un mare conteso}

La Sicilia non era mai stata semplicemente «un'isola». Dal IX all'XI
secolo era stata un emirato islamico di straordinaria complessità
culturale. Con la conquista normanna (1061–1091) era diventata il centro
di un regno che teneva insieme culture araba, greca, latina e normanna
in un equilibrio precario ma creativo. Federico II di Svevia (1194–1250)
aveva fatto di questo regno un laboratorio politico e intellettuale unico
in Europa, trasformando Palermo in una delle capitali più sofisticate del
mondo medievale.

La morte di Federico II lasciò un vuoto. I suoi eredi — Corrado IV, poi
il giovane Corradino — non riuscirono a tenere il potere. Papa Innocenzo
IV, che aveva scomunicato Federico II e combattuto gli Svevi per decenni,
non aveva alcuna intenzione di permettere che il regno di Sicilia restasse
in mano a quella dinastia. La soluzione papale fu radicale: offrire il
regno a una casa reale straniera, abbastanza potente da sconfiggere gli
Svevi ma abbastanza dipendente da Roma da restare controllabile.

La scelta cadde su Carlo d'Angiò, fratello di Luigi IX di Francia.

\secsep

\section{Carlo d'Angiò: un re impresario}

Carlo d'Angiò non era semplicemente un conquistatore. Era qualcosa di più
moderno e più pericoloso: un impresario politico con visioni imperiali.
Aveva sconfitto Manfredi di Svevia nella battaglia di Benevento (1266) e
fatto decapitare il giovanissimo Corradino nella piazza del Mercato di
Napoli (1268), chiudendo la dinastia sveva con una crudeltà calcolata che
fece scalpore in tutta Europa.

Installatosi come re di Sicilia e di Napoli, il suo progetto era la
costruzione di un empire méditerranéen con proiezioni verso Costantinopoli
e il Levante. Nel 1267 aveva firmato il Trattato di Viterbo con Baldovino
II, l'imperatore latino di Costantinopoli in esilio, ottenendo il diritto
di restaurare l'Impero Latino d'Oriente. Nel 1270 aveva partecipato —
e praticamente sabotato — l'ultima crociata di suo fratello Luigi IX,
morto di peste davanti a Tunisi. Nel 1272 era diventato re d'Albania.
Nel 1277 aveva acquistato il titolo di re di Gerusalemme.

\citastorica{%
  Carlo aveva fatto del Mediterraneo una questione personale.
  La Sicilia era la sua base logistica, il suo granaio, il suo arsenale.
  I siciliani erano i suoi strumenti.%
}{Steven Runciman, \textit{The Sicilian Vespers}, 1958}

Il problema era che i siciliani — baroni, clero, artigiani, contadini —
non avevano alcun interesse ad essere strumenti di nessuno.

\subsection{Il peso angioino sulla Sicilia}

Il dominio angioino si tradusse in un sistema di prelievo fiscale e di
requisizioni militari che non aveva precedenti nella storia dell'isola.
Per finanziare le sue campagne, Carlo impose tasse straordinarie con una
frequenza che rese impossibile qualsiasi pianificazione economica per
la popolazione locale. I funzionari francesi — i cosiddetti \textit{provenzali}
o \textit{francesi}, termine che i siciliani usavano indistintamente per
tutti i continentali al servizio di Carlo — erano percepiti come corpi
estranei ostili, spesso arroganti, talvolta violenti.

\nota{Per il GM — Fazione Popolo}{
  Questi meccanismi di sopruso quotidiano sono la materia prima della
  fazione Popolo. I personaggi di questo tavolo non hanno grandi ideologie:
  hanno conti da pagare, donne e figlie da proteggere, pesi da portare.
  Il momento del Vespro non nasce da un piano — nasce da un'umiliazione
  troppa.
}

Le requisizioni di navi, cavalli e vettovaglie per la spedizione contro
Bisanzio — prevista per la primavera del 1282 — avevano raggiunto livelli
insostenibili. La spedizione era il progetto militare più ambizioso che
Carlo avesse mai concepito: una flotta di 300 navi, un esercito di oltre
15.000 uomini. Tutto questo pesava sulla Sicilia.

\secsep

\section{La cospirazione: realtà, leggenda, strumento narrativo}

\subsection{Giovanni da Procida e la rete aragonese}

Giovanni da Procida (1210 ca.–1298) è il protagonista della versione
«cospirazionista» dei Vespri. Medico, giurista, diplomatico, era stato
uno dei più fedeli consiglieri di Federico II e poi di Manfredi. Dopo
la caduta degli Svevi aveva trovato rifugio alla corte di Pietro III
d'Aragona, che aveva sposato Costanza, figlia di Manfredi, e riteneva
di avere diritti legittimi sul regno di Sicilia.

Secondo la tradizione storiografica — cristallizzata nella \textit{Chronica
de rebus in Apulia gestis} di Bartolomeo di Neocastro e nella
\textit{Leggenda di messer Gianni di Procita} di Giovanni Villani — da
Procida avrebbe tessuto una rete di contatti che collegava la corte
aragonese di Barcellona, il papa Niccolò III (ostile a Carlo), l'imperatore
bizantino Michele VIII Paleologo (che aveva tutto l'interesse a bloccare
la spedizione di Carlo contro Costantinopoli) e i baroni siciliani.

\nota{Dibattito storiografico — usare nel gioco}{
  La storiografia moderna è divisa. Runciman accetta l'esistenza di una
  cospirazione strutturata ma sottolinea che i congiurati non controllavano
  i tempi: la rivolta del 30 marzo scoppiò prima del previsto, innescata
  da un incidente che nessuno aveva pianificato. Abulafia tende a ridurre
  il peso della cospirazione e ad amplificare quello del risentimento
  popolare spontaneo. \textbf{Nel gioco, questa ambiguità è una risorsa}:
  la fazione Congiurati ha un piano, ma il piano non è perfetto, e gli
  altri tavoli agiranno con o senza di loro.
}

\subsection{Michele VIII Paleologo e l'oro bizantino}

L'imperatore bizantino Michele VIII Paleologo aveva riconquistato
Costantinopoli nel 1261, ponendo fine all'Impero Latino d'Oriente. La
prospettiva di vedersi di nuovo strappare la città da un esercito angioino
era per lui un incubo esistenziale. Le fonti indicano che Michele finanziò
attivamente la rete di da Procida, inviando oro e lettere criptate attraverso
intermediari.

Quanto fosse coordinato questo sostegno — se si trattasse di un'operazione
di intelligence pianificata o di un finanziamento opportunistico a forze
già esistenti — rimane oggetto di dibattito. Per i giocatori del tavolo
Congiurati, l'oro di Michele è reale. La domanda è quanto ci si possa
fidare di chi lo manda.

\subsection{Il ruolo della Chiesa}

Il papato del 1282 non era monolitico. Papa Martino IV (eletto nel 1281)
era un francese, creatura di Carlo, che aveva scomunicato Michele VIII
per ragioni politiche più che teologiche. Ma la curia romana aveva al
suo interno correnti ostili all'egemonia angioina: vescovi siciliani,
cardinali italiani, frati che avevano visto le loro comunità schiacciate
dalle requisizioni.

Il clero siciliano era in una posizione di straordinaria ambiguità. Le
chiese erano luoghi di rifugio e di comunicazione clandestina. I frati
predicatori e minori circolavano liberamente tra le città. I vescovi
avevano rapporti diretti con Roma che bypassavano l'amministrazione
angioina. Erano, in potenza, una rete di intelligence già esistente —
ma i suoi membri avevano obbedienza a Roma, non ad Aragona, e i loro
interessi erano più sottili e più longevi di quelli dei congiurati.

\secsep

\section{Palermo nel febbraio 1282}

\subsection{La città e i suoi spazi}

Palermo nel 1282 era una città di circa 50.000 abitanti, una delle più
grandi d'Europa. La sua planimetria portava i segni sovrapposti di sette
secoli di dominazioni diverse: la città araba con i suoi \textit{hara}
(quartieri), la città normanna con le sue chiese dorate, i sobborghi
cresciuti intorno alle porte. Il porto era il suo cuore pulsante.

\begin{multicols}{2}
\noindent\textbf{Quartieri principali:}
\begin{itemize}[leftmargin=*, nosep]
  \item \textit{Cassaro} — il centro storico, sede del palazzo reale
  \item \textit{Kalsa} — il quartiere arabo, vicino al porto
  \item \textit{Seralcadio} — artigiani e mercanti
  \item \textit{Albergheria} — il quartiere povero, più densamente popolato
\end{itemize}

\columnbreak

\noindent\textbf{Luoghi chiave:}
\begin{itemize}[leftmargin=*, nosep]
  \item \textit{Chiesa di Santo Spirito} — scintilla del Vespro
  \item \textit{Palazzo dei Normanni} — sede del potere angioino
  \item \textit{Porto di Palermo} — flotta di Carlo in allestimento
  \item \textit{Cattedrale} — sede dell'arcivescovo
\end{itemize}
\end{multicols}

\subsection{La tensione nell'aria}

Le settimane precedenti il 30 marzo erano caratterizzate da una tensione
crescente che chiunque vivesse a Palermo poteva sentire, anche senza
sapere nulla di Giovanni da Procida o di Pietro d'Aragona. I soldati
francesi erano diventati più aggressivi nelle perquisizioni. Le voci di
nuove requisizioni per la spedizione di primavera circolavano nei mercati.
Qualcuno diceva che Carlo era in procinto di imporre un nuovo censimento
— il che significava nuove tasse.

I funzionari angioini che presidiavano la città avevano capito che qualcosa
stava covando sotto la superficie. Jean de Saint-Rémy, il giustiziere
regio di Palermo, aveva intensificato i pattugliamenti e aumentato le
guardie alle porte della città. Ma non sapeva dove guardare, perché
non c'era un unico centro della cospirazione — o forse non c'era nessuna
cospirazione, solo un popolo che non ne poteva più.

\secsep

\section{Il 30 marzo 1282: l'inizio della fine}

\subsection{I Vespri alla chiesa di Santo Spirito}

La sera del lunedì di Pasqua — 30 marzo 1282 — i palermitani si riunirono
alla chiesa di Santo Spirito (detta anche chiesa dei Vespri) per la
processione festiva. Era una delle rare occasioni in cui una folla numerosa
si raccoglieva fuori dalle mura della città, in un contesto semi-festivo,
con le guardie angioine presenti per il controllo d'ordine.

Ciò che accadde esattamente — chi iniziò, se ci fu una provocazione
deliberata o un incidente spontaneo — è proprio il tipo di questione
storica che non può essere risolta con certezza. Le fonti concordano
sull'episodio della perquisizione delle donne da parte dei soldati
francesi, interpretata come violenza e oltraggio; concordano sull'esplosione
improvvisa della violenza; concordano sul fatto che in poche ore la
rivolta si era estesa all'intera città.

\citastorica{%
  Tutta la gente della città, come se avesse aspettato solo quel segnale,
  si levò contro i Francesi con coltelli e armi di fortuna, gridando:
  Mora, mora li Francisi.%
}{Giovanni Villani, \textit{Nuova Cronica}, XIV sec.}

Entro la mattina del 31 marzo, quasi tutti i francesi a Palermo — circa
2.000 persone, soldati, funzionari e civili — erano stati uccisi. La
rivolta si diffuse rapidamente nelle altre città siciliane. Il progetto
imperiale di Carlo d'Angiò era finito.

\subsection{Dopo i Vespri: la guerra del Vespro}

Le conseguenze della rivolta si protrassero per decenni. Pietro III
d'Aragona sbarcò in Sicilia nell'agosto del 1282, rivendicando il trono
in nome di sua moglie Costanza. Iniziò la cosiddetta \textit{Guerra del
Vespro} (1282–1302), uno dei conflitti più costosi e prolungati del
Mediterraneo medievale. Carlo d'Angiò morì nel 1285, con il suo sogno
imperiale in frantumi. La Sicilia rimase aragonese; il regno di Napoli
restò angioino. L'isola era diventata una pedina in un gioco più grande
di lei — ma questa volta aveva scelto da che parte stare.

\nota{Nota di design — perché fermarsi al 30 marzo}{
  \textit{Vespri 1282} si concentra sulla \textbf{vigilia} della rivolta,
  non sulla rivolta stessa. Il 30 marzo è il momento di climax, ma la
  storia che interessa ai giocatori è quella che lo precede: le scelte,
  le alleanze, i tradimenti, i piani che funzionano o crollano. La violenza
  della notte del Vespro è fuori scena — è ciò verso cui tutto tende,
  non ciò che si gioca.
}

\filettodoppio

